\section{Half Term}

\paragraph{Overview:}
This half term I was fortunate enough to spend one week with the team at Cambridge Kinetics, a team of graduates who have built a tool \href{https://www.kinabase.com}{Kinabase} which leverages AI alongside forward thinking data handling to provide companies both large and small the ability to take control of their data and make it work better for them, with these learnings In tow I have chosen to produce an artifact project and to build some of my own software influenced by my learnings of human psychology

\paragraph{Achievements, Thoughts and Ideas:}
Prior to Kinabase, software enterprises relied on large, expensive and specialised tools to manage their business data. These tools also required heavy use of human oversight both to setup and deploy. I joined the team as they have recently received a grant from Innovate UK to further develop their AI tooling and features. This was a phenomenal opportunity which provided me so much insight. I was able to speak to engineers and management who shared with me their learnings about the place of AI within industry. This first hand insight  has left me with a large reading list and numerous spaces I can follow as the technology continually progresses

Following my work experience with an AI focused team such as Cambridge kinetics, I began to approach my other experiences with a view as to what else could AI do? For the past 6+ years I have been volunteering at theatres around me, mixing sound for vocals, orchestras and helping with every aspect of technical backstage productions. This provided another unique opportunity to research the impact AI could have. Theatre being a traditionally human, artistic and emotional industry rarely finds opportunities to integrate analytical tools such as ai into its workflow. I am equally fortunate that my mentor and manager within this show is by trade an engineer and was able to help me apply my ideas and quickly written software programs into functional tools that now form a core part of our workflow. Whilst AI wasn’t able to replace my role as a sound engineer during shows, It excelled as an aid, highlighting show segments which it estimated to require careful human intervention to manage, and automatically generating mic change, organisation templates


During and following my work experience I have completed 4 further pieces of research on the topic of AI (\nameref{research:AIReading}, \nameref{resarch:NaturalLanguage}, \nameref{research:HuggingFace}, \nameref{research:AIProject})

In addition to learning about the modern world of computing, I also had the opportunity to speak to the team about key factors they focus on when looking at how users utilise their products, this made the remainder of my research change strategy slightly and I began researching how humans process the same data as machines but rather than using databases, using notes and records (\nameref{research:NotetakingMethods}, \nameref{research:NoteSearchingMethods}) This provided some fascinating insight and I really want to integrate this human psychology field into my EP

\paragraph{What went well?}
This entire half term was very successful and I am so excited to continue my EPQ and build a tool I am truly proud of
\paragraph{What did I learn?}
One of my biggest takeaways this week, was use what you make. Only by truly understanding and relying on your software do you spot specific bugs, issues and areas for improvement / development. Following this philosophy, once I have my MVP (minimum viable product) I will use both this OneNote notebook and my software in parallel to track all of my sources, notes and related materials. The reason I will use both is that I understand my code will not be perfect, It is likely at some point I will find a large enough error the current software iteration cant overcome, therefore this notebook is the reliable and tested backup which I can rely on. By the end of this project this notebook should be made redundant as all information can be more easily searched, stored, presented and submitted through my software
\paragraph{Conclusions:}
Over the half term I reapproached what I wanted to achieve from my EPQ project. While my original idea of a dissertation surrounding the impact of Ai in our everyday lives still interests me, my work experience and related research made me want to explore in much greater depth what these tools could be, right now and in the future. For this reason I have decided to really stretch myself and produce my own software which integrates my findings. Therefore my EPQ Is going to be an artifact project exploring the role of technology within human learning and knowledge acquisition

\paragraph{Plans for next week:}
Having spent the half term learning all about software development,  AI and user interaction I now need to decide on what my final software project will be
